\documentclass[12pt,a4paper]{article}

\usepackage{cmap}
\usepackage[T2A]{fontenc}
\usepackage[utf8]{inputenc}
\usepackage[russian]{babel}
\usepackage{amsmath,amssymb}
\usepackage[a4paper,margin=2cm]{geometry}

\newcommand{\R}{\mathbb{R}}
\newcommand{\eps}{\varepsilon}
\newcommand{\dd}{\,d}

\setlength{\parindent}{0pt}
\setlength{\parskip}{0.45em}
\sloppy

\begin{document}

\begin{center}
{\Large\bfseries Билет 11}
\end{center}

\noindent\fbox{%
  \begin{minipage}{\dimexpr\linewidth-2\fboxsep-2\fboxrule\relax}
  \normalfont
Равномерная сходимость функциональных последовательностей и рядов.
Критерий Коши равномерной сходимости. Непрерывность суммы равномерно
сходящегося ряда из непрерывных функций. Интегрирование и дифференцирование
функциональных последовательностей и рядов. Признаки Вейерштрасса, Дирихле и
Абеля равномерной сходимости функциональных рядов.
  \end{minipage}%
}

\section*{Краткий план}

Сначала задаются поточечная и равномерная сходимость функциональной
последовательности и функционального ряда. Затем формулируется критерий Коши:
для последовательности он контролирует разности \(f_{n+p}-f_n\), а для ряда --
хвостовые суммы \(\sum_{k=n+1}^{n+p}u_k\), причем оценки должны быть
равномерными по \(x\). После этого доказываются свойства равномерных пределов:
непрерывность, переход к пределу под знаком интеграла и достаточное условие
почленного дифференцирования. Для рядов эти свойства применяются к частичным
суммам. В конце идут признаки равномерной сходимости рядов: Вейерштрасса,
Дирихле и Абеля.

\section*{Определения}

\textbf{Поточечная и равномерная сходимость последовательности.}
Пусть на множестве \(X\) заданы функции \(f_n:X\to\R\) или
\(f_n:X\to\mathbb{C}\).
Последовательность \(\{f_n\}\) сходится к \(f\) поточечно на \(X\), если
\[
  \forall x\in X\quad \lim_{n\to\infty}f_n(x)=f(x).
\]
Она сходится к \(f\) равномерно на \(X\), если
\[
  \forall\eps>0\ \exists N\in\mathbb N:\ 
  \forall n\ge N\ \forall x\in X\quad |f_n(x)-f(x)|\le \eps.
\]
Эквивалентно,
\[
  \sup_{x\in X}|f_n(x)-f(x)|\to0,\qquad n\to\infty.
\]

\textbf{Функциональный ряд.}
Для функций \(u_k:X\to\R\) или \(u_k:X\to\mathbb{C}\) ряд
\[
  \sum_{k=1}^{\infty}u_k(x)
\]
называется функциональным. Его частичные суммы:
\[
  S_n(x)=\sum_{k=1}^n u_k(x).
\]
Ряд сходится равномерно на \(X\), если последовательность \(\{S_n\}\)
равномерно сходится на \(X\) к сумме ряда \(S(x)\). Если ряд поточечно
сходится, его остаток:
\[
  r_n(x)=S(x)-S_n(x)=\sum_{k=n+1}^{\infty}u_k(x).
\]
Тогда ряд сходится равномерно тогда и только тогда, когда
\[
  \sup_{x\in X}|r_n(x)|\to0.
\]

\textbf{Равномерная ограниченность.}
Последовательность функций \(\{g_n\}\) равномерно ограничена на \(X\), если
существует \(C>0\), такое что
\[
  |g_n(x)|\le C\qquad \forall n\in\mathbb N,\ \forall x\in X.
\]

\section*{Теоремы}

\textbf{1. Критерий Коши для функциональной последовательности.}
Последовательность \(\{f_n\}\) равномерно сходится на \(X\) тогда и только
тогда, когда
\[
  \forall\eps>0\ \exists N\in\mathbb N:\ 
  \forall n\ge N\ \forall p\in\mathbb N\ \forall x\in X
  \quad |f_{n+p}(x)-f_n(x)|\le \eps.
\]

\textbf{2. Критерий Коши для функционального ряда.}
Ряд \(\sum_{k=1}^{\infty}u_k(x)\) равномерно сходится на \(X\) тогда и только
тогда, когда
\[
  \forall\eps>0\ \exists N\in\mathbb N:\ 
  \forall n\ge N\ \forall p\in\mathbb N\ \forall x\in X
  \quad
  \left|\sum_{k=n+1}^{n+p}u_k(x)\right|\le \eps.
\]
В частности, из равномерной сходимости ряда следует необходимое условие
\[
  u_n(x)\rightrightarrows 0\quad \text{на }X.
\]

\textbf{3. Непрерывность равномерного предела и суммы ряда.}
Если \(f_n\) непрерывны на \(X\) и \(f_n\rightrightarrows f\) на \(X\), то
\(f\) непрерывна на \(X\). Поэтому если все \(u_k\) непрерывны на \(X\), а ряд
\(\sum u_k(x)\) равномерно сходится на \(X\), то его сумма \(S(x)\) непрерывна
на \(X\).

\textbf{4. Интегрирование функциональных последовательностей и рядов.}
Пусть \(f_n\) непрерывны на \([a,b]\) и \(f_n\rightrightarrows f\) на
\([a,b]\). Тогда
\[
  \lim_{n\to\infty}\int_a^b f_n(x)\dd x
  =
  \int_a^b f(x)\dd x.
\]
Если \(u_k\) непрерывны на \([a,b]\) и ряд \(\sum u_k(x)\) равномерно
сходится на \([a,b]\), то его можно интегрировать почленно:
\[
  \int_a^b\left(\sum_{k=1}^{\infty}u_k(x)\right)\dd x
  =
  \sum_{k=1}^{\infty}\int_a^b u_k(x)\dd x.
\]

\textbf{5. Дифференцирование функциональных последовательностей и рядов.}
Пусть \(f_n\in C^1[a,b]\), последовательность чисел \(f_n(x_0)\) сходится хотя
бы в одной точке \(x_0\in[a,b]\), а производные \(f_n'\) равномерно сходятся
на \([a,b]\) к функции \(\varphi\). Тогда \(f_n\) равномерно сходится на
\([a,b]\) к функции
\[
  f(x)=A+\int_{x_0}^x\varphi(t)\dd t,\qquad
  A=\lim_{n\to\infty}f_n(x_0),
\]
причем \(f\in C^1[a,b]\) и
\[
  f'(x)=\varphi(x)=\lim_{n\to\infty}f_n'(x).
\]
Для рядов: если \(u_k\in C^1[a,b]\), ряд \(\sum u_k(x_0)\) сходится хотя бы в
одной точке \(x_0\), а ряд производных \(\sum u_k'(x)\) равномерно сходится
на \([a,b]\), то ряд \(\sum u_k(x)\) равномерно сходится на \([a,b]\) и
\[
  \left(\sum_{k=1}^{\infty}u_k(x)\right)'
  =
  \sum_{k=1}^{\infty}u_k'(x).
\]

\textbf{6. Признак Вейерштрасса.}
Если
\[
  |u_k(x)|\le a_k\qquad \forall k\in\mathbb N,\ \forall x\in X
\]
и числовой ряд \(\sum_{k=1}^{\infty}a_k\) сходится, то функциональный ряд
\(\sum_{k=1}^{\infty}u_k(x)\) сходится равномерно на \(X\).

\textbf{7. Признак Дирихле.}
Пусть на \(X\) заданы последовательности \(a_k(x)\) и \(b_k(x)\). Если
\[
  A_n(x)=\sum_{k=1}^n a_k(x)
\]
равномерно ограничены на \(X\), то есть \(|A_n(x)|\le C\), если
\[
  b_k(x)\rightrightarrows 0\quad \text{на }X,
\]
и при каждом фиксированном \(x\in X\) последовательность \(b_k(x)\)
монотонно убывает:
\[
  b_{k+1}(x)\le b_k(x),
\]
то ряд
\[
  \sum_{k=1}^{\infty}a_k(x)b_k(x)
\]
сходится равномерно на \(X\).

\textbf{8. Признак Абеля.}
Пусть ряд \(\sum_{k=1}^{\infty}a_k(x)\) равномерно сходится на \(X\),
последовательность \(b_k(x)\) равномерно ограничена:
\[
  |b_k(x)|\le C,
\]
и при каждом \(x\in X\) она монотонна:
\[
  b_{k+1}(x)\le b_k(x).
\]
Тогда ряд
\[
  \sum_{k=1}^{\infty}a_k(x)b_k(x)
\]
сходится равномерно на \(X\).

\section*{Доказательства}

\textbf{Критерий Коши для последовательности.}
Пусть \(f_n\rightrightarrows f\). Для заданного \(\eps>0\) выберем \(N\) так,
что при \(m\ge N\)
\[
  |f_m(x)-f(x)|\le \frac{\eps}{2}\qquad \forall x\in X.
\]
Тогда при \(n\ge N\) и \(p\in\mathbb N\)
\[
  |f_{n+p}(x)-f_n(x)|
  \le |f_{n+p}(x)-f(x)|+|f_n(x)-f(x)|
  \le \eps
\]
равномерно по \(x\). Значит, условие Коши выполнено.

Обратно, пусть выполнено условие Коши. Тогда для каждого фиксированного
\(x\in X\) числовая последовательность \(\{f_n(x)\}\) фундаментальна, значит
имеет предел; обозначим его \(f(x)\). В условии Коши фиксируем \(n\ge N\) и
\(x\), а затем переходим к пределу при \(p\to\infty\):
\[
  |f_n(x)-f(x)|\le \eps.
\]
Так как \(N\) не зависит от \(x\), получаем \(f_n\rightrightarrows f\) на
\(X\).

\textbf{Критерий Коши для ряда.}
Применяем предыдущий критерий к частичным суммам
\[
  S_n(x)=\sum_{k=1}^n u_k(x).
\]
Разность частичных сумм равна
\[
  S_{n+p}(x)-S_n(x)=\sum_{k=n+1}^{n+p}u_k(x),
\]
поэтому условие Коши для \(\{S_n\}\) в точности совпадает с указанным условием
для хвостовых сумм ряда.

\textbf{Непрерывность равномерного предела.}
Зафиксируем \(x_0\in X\) и \(\eps>0\). По равномерной сходимости выберем
\(N\), такое что
\[
  |f(x)-f_N(x)|<\frac{\eps}{3}\qquad \forall x\in X.
\]
Так как \(f_N\) непрерывна в \(x_0\), существует \(\delta>0\), для которого
при \(x\in X\), \(|x-x_0|<\delta\),
\[
  |f_N(x)-f_N(x_0)|<\frac{\eps}{3}.
\]
Тогда
\[
  |f(x)-f(x_0)|
  \le |f(x)-f_N(x)|+|f_N(x)-f_N(x_0)|+|f_N(x_0)-f(x_0)|
  <\eps.
\]
Значит, \(f\) непрерывна. Для ряда частичные суммы \(S_n\) непрерывны как
конечные суммы непрерывных функций, а \(S_n\rightrightarrows S\); поэтому
сумма \(S\) непрерывна.

\textbf{Интегрирование.}
Из уже доказанной теоремы о непрерывности предельной функции следует, что
\(f\) непрерывна на \([a,b]\), значит интегрируема. Кроме того,
\[
  \left|\int_a^b f_n(x)\dd x-\int_a^b f(x)\dd x\right|
  \le
  \int_a^b |f_n(x)-f(x)|\dd x
  \le
  (b-a)\sup_{x\in[a,b]}|f_n(x)-f(x)|\to0.
\]
Для ряда применяем это к частичным суммам \(S_n=\sum_{k=1}^n u_k\):
\[
  \int_a^b S(x)\dd x
  =
  \lim_{n\to\infty}\int_a^b S_n(x)\dd x
  =
  \lim_{n\to\infty}\sum_{k=1}^n\int_a^b u_k(x)\dd x.
\]

\textbf{Дифференцирование.}
Пусть \(f_n'\rightrightarrows\varphi\) на \([a,b]\) и
\(f_n(x_0)\to A\). Так как \(f_n'\) непрерывны, по теореме о непрерывности
равномерного предела \(\varphi\) непрерывна. Определим
\[
  f(x)=A+\int_{x_0}^x\varphi(t)\dd t.
\]
По формуле Ньютона--Лейбница
\[
  f_n(x)=f_n(x_0)+\int_{x_0}^x f_n'(t)\dd t.
\]
Следовательно,
\[
  |f_n(x)-f(x)|
  \le
  |f_n(x_0)-A|+
  \left|\int_{x_0}^x (f_n'(t)-\varphi(t))\dd t\right|
  \le
  |f_n(x_0)-A|+(b-a)\sup_{t\in[a,b]}|f_n'(t)-\varphi(t)|.
\]
Правая часть стремится к нулю и не зависит от \(x\), поэтому
\(f_n\rightrightarrows f\). Из определения \(f\) и непрерывности \(\varphi\)
следует \(f'(x)=\varphi(x)\).

Для ряда берутся частичные суммы \(S_n=\sum_{k=1}^n u_k\). Условие сходимости
\(\sum u_k(x_0)\) дает сходимость \(S_n(x_0)\), а равномерная сходимость
\(\sum u_k'\) дает равномерную сходимость \(S_n'=\sum_{k=1}^n u_k'\). По
только что доказанной теореме \(S_n\) равномерно сходится и
\[
  \left(\lim_{n\to\infty}S_n(x)\right)'
  =
  \lim_{n\to\infty}S_n'(x),
\]
то есть ряд можно дифференцировать почленно.

\textbf{Признак Вейерштрасса.}
По критерию Коши для числового ряда \(\sum a_k\):
\[
  \forall\eps>0\ \exists N:\ \forall n\ge N\ \forall p\in\mathbb N
  \quad \sum_{k=n+1}^{n+p}a_k\le \eps.
\]
Тогда при всех \(x\in X\)
\[
  \left|\sum_{k=n+1}^{n+p}u_k(x)\right|
  \le
  \sum_{k=n+1}^{n+p}|u_k(x)|
  \le
  \sum_{k=n+1}^{n+p}a_k
  \le \eps.
\]
По критерию Коши для функциональных рядов \(\sum u_k(x)\) сходится
равномерно.

\textbf{Признак Дирихле.}
Положим
\[
  A_m(x)=\sum_{k=1}^m a_k(x),\qquad |A_m(x)|\le C.
\]
Для хвоста ряда введем
\[
  B_j(x)=\sum_{k=n+1}^{j}a_k(x)=A_j(x)-A_n(x),\qquad n+1\le j\le n+p.
\]
Тогда \(|B_j(x)|\le 2C\). По преобразованию Абеля
\[
  \sum_{k=n+1}^{n+p}a_k(x)b_k(x)
  =
  B_{n+p}(x)b_{n+p}(x)
  +
  \sum_{k=n+1}^{n+p-1}B_k(x)\bigl(b_k(x)-b_{k+1}(x)\bigr).
\]
Так как \(b_k(x)\downarrow 0\), то \(b_k(x)\ge0\), и поэтому
\[
  \left|\sum_{k=n+1}^{n+p}a_k(x)b_k(x)\right|
  \le
  2C b_{n+p}(x)+2C\sum_{k=n+1}^{n+p-1}(b_k(x)-b_{k+1}(x))
  \le
  2C b_{n+1}(x).
\]
Из \(b_n\rightrightarrows0\) следует, что правая часть стремится к нулю
равномерно по \(x\). Значит, выполнен критерий Коши для ряда
\(\sum a_k(x)b_k(x)\).

\textbf{Признак Абеля.}
Пусть
\[
  R_n(x)=\sum_{k=n+1}^{\infty}a_k(x).
\]
Из равномерной сходимости \(\sum a_k(x)\) следует
\[
  M_n:=\sup_{j\ge n,\ x\in X}|R_j(x)|\to0.
\]
Для конечного хвоста произведенного ряда снова применяем преобразование
Абеля. Так как \(|b_k(x)|\le C\) и монотонная последовательность \(b_k(x)\)
имеет на любом отрезке индексов суммарное изменение не больше \(2C\), получаем
оценку вида
\[
  \left|\sum_{k=n+1}^{n+p}a_k(x)b_k(x)\right|
  \le
  4C M_n.
\]
Правая часть не зависит от \(x\) и стремится к нулю. Следовательно, по
критерию Коши ряд \(\sum a_k(x)b_k(x)\) сходится равномерно на \(X\).

\section*{Финальная устная версия}

Равномерная сходимость последовательности \(f_n\to f\) на \(X\) означает, что
номер \(N\) в определении предела можно выбрать один для всех \(x\in X\):
\[
  \forall\eps>0\ \exists N:\ n\ge N,\ x\in X
  \Rightarrow |f_n(x)-f(x)|\le\eps.
\]
Эквивалентно, \(\sup_X|f_n-f|\to0\). Функциональный ряд
\(\sum u_k(x)\) сходится равномерно, если равномерно сходится
последовательность его частичных сумм \(S_n=\sum_{k=1}^n u_k\). Для поточечно
сходящегося ряда это равносильно \(\sup_X|r_n|\to0\), где
\(r_n=\sum_{k=n+1}^{\infty}u_k\).

Критерий Коши для последовательностей:
\[
  \{f_n\}\text{ равномерно сходится}
  \Longleftrightarrow
  \forall\eps>0\ \exists N:\ n\ge N,\ p\in\mathbb N,\ x\in X
  \Rightarrow |f_{n+p}(x)-f_n(x)|\le\eps.
\]
Доказательство такое же, как для числовых последовательностей, но оценки
держатся равномерно по \(x\): из равномерной сходимости разность оценивается
через \(f\), а из условия Коши сначала получается поточечный предел
\(f(x)\), затем переходом \(p\to\infty\) выходит равномерная оценка
\(|f_n(x)-f(x)|\le\eps\). Для ряда критерий Коши записывается для хвостов:
\[
  \left|\sum_{k=n+1}^{n+p}u_k(x)\right|\le\eps
\]
при всех больших \(n\), всех \(p\) и всех \(x\).

Если \(f_n\) непрерывны и \(f_n\rightrightarrows f\), то \(f\) непрерывна:
\[
  |f(x)-f(x_0)|\le |f-f_N|(x)+|f_N(x)-f_N(x_0)|+|f_N-f|(x_0),
\]
и три слагаемых делаются малыми. Поэтому равномерно сходящийся ряд из
непрерывных функций имеет непрерывную сумму.

На отрезке \([a,b]\) равномерный предел непрерывных функций можно
интегрировать предельным переходом:
\[
  \left|\int_a^b f_n-\int_a^b f\right|
  \le (b-a)\sup_{[a,b]}|f_n-f|\to0.
\]
Отсюда для равномерно сходящегося ряда непрерывных функций следует
\[
  \int_a^b\sum_{k=1}^{\infty}u_k(x)\dd x
  =
  \sum_{k=1}^{\infty}\int_a^b u_k(x)\dd x.
\]

Для дифференцирования нужны более сильные условия: если \(f_n\in C^1[a,b]\),
\(f_n(x_0)\) сходится хотя бы в одной точке, а \(f_n'\rightrightarrows\varphi\),
то
\[
  f_n(x)=f_n(x_0)+\int_{x_0}^x f_n'(t)\dd t
\]
показывает равномерную сходимость \(f_n\) к
\[
  f(x)=A+\int_{x_0}^x\varphi(t)\dd t.
\]
Следовательно, \(f'(x)=\varphi(x)=\lim f_n'(x)\). Для ряда это дает правило:
если \(\sum u_k(x_0)\) сходится в одной точке, а \(\sum u_k'\) равномерно
сходится, то \(\sum u_k\) равномерно сходится и
\[
  \left(\sum u_k\right)'=\sum u_k'.
\]

Признак Вейерштрасса: если \(|u_k(x)|\le a_k\) для всех \(x\), а числовой ряд
\(\sum a_k\) сходится, то \(\sum u_k(x)\) сходится равномерно. Доказательство:
хвост функционального ряда по модулю не больше хвоста сходящегося числового
ряда.

Признак Дирихле: если частичные суммы \(\sum_{k=1}^n a_k(x)\) равномерно
ограничены, \(b_k(x)\rightrightarrows0\), и \(b_k(x)\) монотонно убывает по
\(k\) при каждом \(x\), то \(\sum a_k(x)b_k(x)\) сходится равномерно. Главное
доказательство -- преобразование Абеля, после которого хвост оценивается
через \(C\sup_X b_{n+1}\to0\).

Признак Абеля: если \(\sum a_k(x)\) равномерно сходится, \(b_k(x)\)
равномерно ограничены и монотонны по \(k\), то \(\sum a_k(x)b_k(x)\) сходится
равномерно. Здесь преобразование Абеля оценивает хвост через равномерный
остаток ряда \(\sum a_k\), который стремится к нулю.

\section*{Источники}

Иванов Г.Е. \emph{Лекции по математическому анализу. Часть 1}:
\texttt{IvGE\_1.pdf}, стр. 281--299, 303.

Основные роли источника: стр. 281--285 -- определения равномерной сходимости
функциональных последовательностей и критерий Коши; стр. 285--288 --
функциональные ряды, остатки, критерий Коши и признак Вейерштрасса; стр.
288--292 -- признаки Дирихле, Лейбница и Абеля; стр. 295--299 --
непрерывность, интегрирование и дифференцирование функциональных
последовательностей и рядов; стр. 303 -- формулировка признака Вейерштрасса
для комплексных функциональных рядов.

Дополнительно был просмотрен \texttt{IvGE\_2.pdf}: прямого материала по
этому билету там не найдено; найденные страницы относятся к интегралам,
зависящим от параметра, и рядам Фурье.

\end{document}
